\title{\LARGE \bf Homolog\'{\i}a de productos semidirectos de grupos.}
\author{\small V. \'{A}lvarez, J.A. Armario y P. Real\\
\small Dpto. Matem\'{a}tica Aplicada I (Universidad de Sevilla)\\
\small Avda. Reina Mercedes, s/n, 41012 Sevilla \\
\small  {\em E-mails:} valvarez@cica.es, armario@cica.es, real@cica.es}
\documentstyle[emlines2,bezier,espanol,12pt]{article}
\date{}
\textheight=615pt
\textwidth=425pt
\oddsidemargin=0pt
\headheight=0pt
\headsep=7pt
\hoffset=0pt
\voffset=0pt
\topmargin=10pt
\marginparsep=10pt
\marginparwidth=30pt
\footheight=10pt
\footskip=40pt
\parskip=.5cm
\setlength{\baselineskip}{8ex}
\newcommand{\ra}{\rightarrow}
\newcommand{\la}{\leftarrow}
\newcommand{\scst}{\scriptscriptstyle}
\newcommand{\ot}{\otimes}
\newcommand{\Z}{{\Bbb Z}}
\newcommand{\Zp}{{\Bbb Z}_p}
\newcommand{\pa}{\partial}
\newcommand{\Fp}{{\Bbb F}_p}
\newcommand{\Fr}{{\Bbb F}_2}
\newcommand{\n}{{\tilde{\mbox{n}}}}
\newtheorem{prp}{Proposici\'{o}n}[section]
\newtheorem{lmm}[prp]{Lema}
\newtheorem{thr}[prp]{Teorema}
\newtheorem{dfn}[prp]{Definici\'{o}n}
\newtheorem{rmk}[prp]{Nota}
\newtheorem{crl}[prp]{Corolario}
\newtheorem{ejemplo}[prp]{Ejemplo}
\newtheorem{notacion}[prp]{Notaci\'{o}n}
\hyphenation{pro-duc-to}
\begin{document}
\bibliographystyle{alpha}
\sloppy
\maketitle
\vspace{0.5cm}

\hyphenation{pro-duc-to}

\begin{abstract}
Tomando como base el estudio sobre la homolog\'{\i}a de fibrados realizado en \cite{LS87}
y la l\'{\i}nea expuesta en \cite{AAR98},
dise$\n$amos aqu\'{\i} un algoritmo para el c\'{a}lculo de la homolog\'{\i}a de ciertos
productos semidirectos de grupos, que comprende tanto grupos discretos como simpliciales.
Completa el trabajo una valoraci\'{o}n de la complejidad que engloba
todo el proceso, anotando circunstancias y condiciones bajo las cuales se
simplificar\'{\i}an los c\'{a}lculos. Como caso particular, analizamos la
homolog\'{\i}a de los grupos dih\'{e}dricos desde esta nueva
perspectiva. Se tiene prevista la implementaci\'{o}n del algoritmo en
el paquete de c\'{a}lculo que provee el {\em Mathematica 3.0}.
\end{abstract}

\begin{center}
{\bf Resumen extendido}
\end{center}

Perseguimos resolver aqu\'{\i} el problema del
c\'{a}lculo de la homolog\'{\i}a de ciertos productos semidirectos
de grupos. Dise$\n$aremos un algoritmo para el c\'{a}lculo de la
homolog\'{\i}a de productos semidirectos $A \rtimes G$ con $G$
abeliano finito y $A$ abeliano no necesariamente finito.

El proceder general se basa en el establecimiento de un isomorfismo que
traslade el problema al \'{a}mbito m\'{a}s combinatorio de la topolog\'{\i}a
simplicial, recurso al que nos vemos con frecuencia abocados los top\'{o}logos
algebraicos: no pasa desapercibido que las contracciones de tipo
Eilenberg-Mac Lane est\'{a}n \'{\i}ntimamente ligadas al c\'{a}lculo
homol\'{o}gico de grupos y provienen del campo simplicial.

Nosotros seguiremos un discurrir an\'{a}logo al descrito en \cite{LS87}. Primero, a partir de
un producto semidirecto de grupos finitos $A \rtimes G$, descompondremos
el grupo simplicial $\bar{W}(A \rtimes G)$ como un producto
cartesiano torcido de la forma
$\bar{W}(A)\times_{\tau}\bar{W}(G)$. A partir de aqu\'{\i},
utilizando t\'{e}cnicas de perturbaci\'{o}n similares a las empleadas
en la referencia anterior, es factible definir una resoluci\'{o}n peque$\n$a
libre para el producto semidirecto de partida. La dificultad que
esconde este en apariencia sencillo resultado estriba en encontrar
la filtraci\'{o}n adecuada que permite que los procesos de
perturbaci\'{o}n utilizados en el transcurso de su demostraci\'{o}n
sean finitos: la filtraci\'{o}n que definimos nosotros difiere
sustancialmente de la utilizada en \cite{LS87}; dado que la
validez de esta \'{u}ltima depend\'{\i}a de unas condiciones
adicionales a verificar por la funci\'{o}n de torsi\'{o}n $\tau$, que
no se dan en el caso particular que nos lleva, del producto
semidirecto de grupos finitos.

Preliminares esenciales en esta trabajo ser\'{a}n los teoremas de
Brown (\cite{Bro59,Bro64}), Eilenberg-Zilber y Eilenberg-Zilber
torcido (\cite{EM54,EZ59}), en la versi\'{o}n que demostrara Shih en
\cite{Shi61}.

\section{Modelos homol\'{o}gicos de productos semidirectos.}

En esta secci\'{o}n nos preocupamos por encontrar un ``modelo
homol\'{o}gico'' para un producto semidirecto de grupos $A \rtimes
G$.

El camino a seguir es el siguiente: primero, concretaremos qu\'{e}
entendemos por modelo homol\'{o}gico de un grupo; en segundo
lugar, definiremos el {\em espacio clasificante geom\'{e}trico},
por medio del cual se traduce la b\'{u}squeda de
estos modelos homol\'{o}gicos a un contexto m\'{a}s adecuado.
Particularizando el estudio para el caso de ciertos productos
semidirectos de grupos, podremos hacer uso de una bater\'{\i}a de
resultados que permitir\'{a}n reducir el problema a un proceso de
perturbaci\'{o}n homol\'{o}gica, del que garantizaremos la parada.

La {\em homolog\'{\i}a} de un grupo suele definirse por medio de
su {\em espacio clasificante}, tambi\'{e}n conocido como {\em
clasificante geom\'{e}trico}. Veamos en qu\'{e} consiste esta
construcci\'{o}n (ver \cite{May67}).

Dado un grupo simplicial $(G,+)$, el {\em espacio clasificante}
$\bar{W}^g(G)$ consiste en el conjunto simplicial definido de la
siguiente manera:
\begin{eqnarray*}
\bar{W}^g_0(G) & = & \{ [ \, ] \} ; \\ \bar{W}^g_n(G) & = &
G_{n-1} \times \cdots \times G_0, \quad n>0; \\ s_0[ \, ] & = &
[e_0]; \\
\partial _i [g_0] & = & [ \, ], \quad i=0,1; \\
\partial _0[g_n, \ldots ,g_0] & = & [g_{n-1}, \ldots ,g_0], \\
\partial _{i+1}[g_n, \ldots ,g_0] & = & [\partial _i g_n, \ldots , \partial
_1 g_{n-i+1},g_{n-i-1} +\partial _0 g_{n-i},g_{n-i-2}, \ldots
,g_0], \\ s_0[g_{n-1}, \ldots ,g_0] & = & [e_n,g_{n-1}, \ldots
,g_0], \\ s_{i+1}[g_n, \ldots ,g_0] & = & [s_i g_n, \ldots ,s_0
g_{n-i},e_{n-i}, g_{n-i-1}, \ldots ,g_0];
\end{eqnarray*}
donde $[\,]$ denota el \'{u}nico elemento de $\bar{W}^g_0(G)$, $e_n$
denota el elemento identidad de $G_n$ y $[g_{n-1}, \ldots , g_0]$
denota un elemento gen\'{e}rico de $\bar{W}^g_n(G)$, para $n>0$.

Por abuso de notaci\'{o}n, se habla del clasificante asociado a un
grupo discreto $G$ como el clasificante $\bar{W}^g(^sG)$ del grupo
simplicial $^sG$ can\'{o}nicamente asociado a $G$, que consiste en
aquel con $^sG_n=G$ para todo $n \geq 0$ y operadores simpliciales
(ambos de cara y de degeneraci\'{o}n) el homomorfismo identidad.

La homolog\'{\i}a de un grupo $G$, ya sea discreto o simplicial,
se puede definir entonces como la homolog\'{\i}a del objeto
$\bar{W}^g(G)$ (ver \cite{Bro82}): es decir, la correspondiente
del complejo de cadenas $C_*(\bar{W}^g(G))$.

La utilizaci\'{o}n de los clasificantes algebraico (\cite{GM74MAMS}) y
geom\'{e}trico (\cite{May67}) y la contracci\'{o}n semicompleta
$C_{\scst \bar{W}-\bar{B}}$ relacionando
clasificante algebraico de un \'{a}lgebra simplicial y construcci\'{o}n Bar
del \'{a}lgebra diferencial graduada correspondiente
(\cite{Rea95,RS94,AAR99}), permiten el dise$\n$o de una maquinaria de
c\'{a}lculo de la homolog\'{\i}a de ciertos grupos, por medio de la
Teor\'{\i}a de Perturbaci\'{o}n Homol\'{o}gica (\cite{GL89,GLS90,
GLS91,Shi61,Bro64,Gug60}).

De hecho, dado un grupo simplicial reducido $G$, el procedimiento
consiste en los siguientes pasos:
\begin{enumerate}
\item Mediante el isomorfismo dado en \cite{GM74MAMS} se
pasa de $C_*(\bar{W}^g(G))$ a $\bar{W}(C_*(G))$.
\item Se aplica la contracci\'{o}n que liga $\bar{W}(C_*(G))$ con
$\bar{B}(C_*(G))$.
\item Se busca alg\'{u}n ``modelo homol\'{o}gico'' del \'{a}lgebra $C_*(G)$
para calcular de la manera m\'{a}s sencilla posible la
homolog\'{\i}a de $\bar{B}C_*(G))$.
\end{enumerate}

En el caso particular de los grupos discretos, el proceder es
sustancialmente distinto, en tanto en cuanto los grupos
simpliciales asociados no son reducidos. Ahora bien, el an\'{a}lisis
detenido de la construcci\'{o}n de la equivalencia de homotop\'{\i}a
$C_{\scst \bar{W}-\bar{B}}$ muestra que en grado simplicial 1 los
morfismos de la contracci\'{o}n se comportan como isomorfismos, en
realidad; luego en el caso de un grupo discreto $G$, resulta que
$C_*(\bar{W}^g(G))$ es can\'{o}nicamente isomorfo a
$\bar{B}(C_*(G))$, por lo que el procedimiento a seguir se reduce
al paso 3 anterior.

Este es el camino que vamos a seguir para establecer la
homolog\'{\i}a de los productos semidirectos de grupos.

La noci\'{o}n de ``modelo homol\'{o}gico'' que consideramos aqu\'{\i}
es una versi\'{o}n m\'{a}s d\'{e}bil de la que definieran Lambe y
Stasheff en \cite{LS87}.

\begin{dfn}
{\em Un {\em modelo homol\'{o}gico} de un grupo (simplicial o
discreto) $G$ es un par $(c,hG)$, donde $c$ es una contracci\'{o}n
que liga $C_*(\bar{W}(G))$ con un DG-m\'{o}dulo $hG$ libre de tipo
finito.}
\end{dfn}

N\'{o}tese que cualesquiera par de modelos son homol\'{o}gicamente
equivalentes, por medio de un puente de homotop\'{\i}a, de base la
construcci\'{o}n bar del complejo de cadenas asociado al grupo en
cuesti\'{o}n. El sentido de calcular un modelo homol\'{o}gico es
precisamente el obtener un complejo que presente una cantidad
peque$\n$a de generadores para as\'{\i} simplificar el proceso
final del c\'{a}lculo de su homolog\'{\i}a.

Todo modelo homol\'{o}gico cuya diferencial sea nula se denomina
{\em modelo minimal}. Claramente, todo modelo minimal es isomorfo
al m\'{o}dulo de homolog\'{\i}a del grupo de partida, de donde la
legitimidad del t\'{e}rmino ``minimal''. De mismo, que se hable de
{\em el modelo minimal}.

Con asiduidad, utilizaremos simplemente el DG-m\'{o}dulo $hG$ para
referirnos a un modelo homol\'{o}gico $(c,hG)$. Otras veces,
adem\'{a}s, consideraremos el complejo $\bar{B}(\Z[G])$ en lugar de
$C_*(\bar{W}(G))$, dado que ambos son isomorfos (\cite{EM53}).

Veamos un modelo homol\'{o}gico para grupos discretos conmutativos,
que ser\'{a} de utilidad a la hora de encontrar un modelo para
productos semidirectos de grupos discretos.

Dado que todo grupo discreto conmutativo factorizar\'{a} como
producto de una ``parte libre'' (cierta potencia $\Z^n$ de $\Z$) y
otra ``de torsi\'{o}n'' (del tipo $\Z_{p_1^{r_1}} \times \cdots
\times \Z_{p_t^{r_t}}$, con cada $p_i$ primo y $r_i \geq 1$),
vamos a centrar nuestro inter\'{e}s en la determinaci\'{o}n de modelos
homol\'{o}gicos para $\bar{B}(\Z [\Z])$ y $\bar{B}(\Z [\Z _{p^r}])$.
A tal efecto, adaptamos a nuestra terminolog\'{\i}a dos resultados
demostrados por Eilenberg y Mac Lane en \cite{EM54}.

Consideremos la colecci\'{o}n de datos $$C_{\scst {\Z}}:\{ \bar{B}
(\Z[{\Z}]),\,\mbox{$E(u,1), \, f_{\scst {\Z}}, \, g_{\scst {\Z}},
\, \phi_{\scst {\Z}}$}\},$$ donde, por extensi\'{o}n lineal, se
define la ``proyecci\'{o}n'' como $$f_{\scst {\Z}}(c_q) = \left\{
\begin{array}{ll} nu, & \mbox{si $q=0$, con $c_0=[n] \in
s( {\Z} )= \bar{B}_0(\Z[\Z]);$}\\ 0, & \mbox{si $q>0$, con $c_q \in
\bar{B}_q(\Z[\Z])$;} \end{array} \right.$$ la ``inyecci\'{o}n'' por
$$g_{\scst {\Z}}(u)=[1];$$ y el ``operador de homotop\'{\i}a''
como $$\phi_{\scst {\Z}}([n_1 \vert \ldots \vert n_k]) = \left\{
\begin{array}{ll} -\sum_{i=1}^{n_1 -1} [1\vert i \vert n_2 \vert
\ldots \vert n_k], & \mbox{ si $n_1
>1$;} \\ 0, & \mbox{ si $n_1 = 1$;} \\ \sum_{i=1}^{\vert n_1
\vert} [1 \vert -i \vert n_2 \vert \ldots \vert n_k], & \mbox{ si
$n_1 < 1$;}
\end{array} \right.$$

Este conjunto de datos define en realidad una contracci\'{o}n. A\'{u}n
m\'{a}s, se verifica el siguiente resultado.

\begin{thr} (\cite{EM54}) \label{5-teorema-bz}
El par $(C_{\scst {\Z}}, E(u,1))$ define un modelo homol\'{o}gico
(de hecho, el minimal), del grupo $\Z$ de los enteros. M\'{a}s
a\'{u}n, la proyecci\'{o}n es un morfismo de DGA-\'{a}lgebras y la
inyecci\'{o}n es un morfismo de DGA-\'{a}lgebras de Hopf (i.e., en
ambos los niveles de \'{a}lgebra y co\'{a}lgebra).
\end{thr}

Ahora queremos hacer lo propio con el complejo $\bar{B}(\Z [\Z
_{p^r}])$. Aunque podr\'{\i}amos considerar como propio modelo
homol\'{o}gico al mismo objeto $\bar{B}(\Z [\Z _{p^r}])$ (se trata
de un DG-m\'{o}dulo libre de tipo finito, a diferencia del caso
anterior, que era libre pero no finitamente generado en cada
grado), se conoce un modelo de magnitud inferior, en el sentido de
que simplifica la complejidad de los procesos computacionales
inmersos en el c\'{a}lculo de su homolog\'{\i}a.

Fijemos un n\'{u}mero primo $p$ y un natural $r$, y tomemos el
conjunto $$C_{\scst {\Z}_{p^r}} : \{ \bar{B}_N ({\Z}[{\Z}_{p^r}])
, \;\; (E(u,1) {\ot} \Gamma(v,2), d_{p^r}), \; f_{\scst
{\Z}_{p^r}}, \; g_{\scst {\Z}_{p^r}}, \; \phi_{\scst
{\Z}_{p^r}}\}.$$ Por extensi\'{o}n lineal, se definen los morfismos
de la siguiente manera: la ``proyecci\'{o}n'', como
$$\begin{array}{rcl} f_{\scst {\Z}_{p^r}}[x_1 \vert y_1 \vert
\ldots \vert x_m \vert y_m] & = & [ \prod_{i=1}^{m}s^2(x_i,y_i)]
\gamma_m(v), \mbox{ }\\ f_{\scst {\Z}_{p^r}}[x_1 \vert y_1 \vert
\ldots \vert x_m \vert y_m \vert z] & = & [ s^1 (z) \prod_{i
=1}^{m} s^2(x_i, y_i)] u \gamma_m (v),
\end{array}$$
donde $s^2: {\Z}_{p^r} \times {\Z}_{p^r} \ra {\Z}$ y
$s^1:{\Z}_{p^r} \ra {\Z}$ son las aplicaciones
\[ s^2(i,j) = \left\{ \begin{array}{ll} 0, & \mbox{  si }\,i+j < p^r; \\
                                        1, & \mbox{  si }\,i+j \geq p^r;\end{array}
\right. \]
\[s^1(i) = i , \;\;\; 0 \leq i<p^r. \]
La ``inyecci\'{o}n'', por $$\begin{array}{lll} \mbox{     } g_{\scst
{\Z}_{p^r}}(u) & = & [1], \mbox{ }\\ g_{\scst {\Z}_{p^r}}(\gamma_k
(v)) & = & \sum_{x_1, \ldots, x_k \in \mbox{\small {\Z}}_{p^r}} [1
\vert x_1 \vert \ldots \vert 1 \vert x_k ],  \\ g_{\scst
{\Z}_{p^r}}(u\gamma_k (v)) & = & \sum_{x_1, \ldots, x_k \in
\mbox{\small {\Z}}_{p^r}}[1 \vert x_1 \vert \ldots \vert 1 \vert
x_k \vert 1 ].
\end{array} $$
Y el ``operador de homotop\'{\i}a'' $$\begin{array}{l} \phi_{\scst
{\Z}_{p^r}} 1 = 0, \;\; \phi_{\scst {\Z}_{p^r}}[x] = -C(x), \;\;
\phi_{\scst {\Z}_{p^r}}[x|y] = -[C(x)|y],
\\ \phi_{\scst {\Z}_{p^r}} [x\vert y \vert \sigma] = -[C(x) \vert y
\vert \sigma] - s^2(x,y) [\sum_{\scst l \in \mbox{\small {\bf
Z}}_{p^r}} [1 \vert l] \vert \phi_{\scst {\Z}_{p^r}}\sigma];
\end{array}$$
donde
\[ C(x) = \sum_{i=1}^{x-1} [1 \vert i]. \]

En estas circunstancias, se verifica este otro teorema.

\begin{thr}(\cite{EM54}) \label{mhp}
El par $(C_{\scst {\Z}_{p^r}},\,(E(u,1) {\ot} \Gamma(v,2),
d_{p^r})) $ es un modelo homol\'{o}gico del grupo c\'{\i}clico
finito ${\Z}_{p^r}$ de orden $p^r$. M\'{a}s a\'{u}n, ambas la
proyecci\'{o}n y la inyecci\'{o}n de la contracci\'{o}n $C_{\scst
{\Z}_{p^r}}$ son morfismos de DGA-\'{a}lgebras.
\end{thr}

La diferencial $d_{p^r}$ es la usual, que act\'{u}a sobre los
generadores $u$ y $v$ en la forma $d_{p^r}(u)=0$ y $d_{p^r}(v)=\pm
p^r u$.

A partir de los resultados anteriores, es inmediato establecer un
modelo homol\'{o}gico de cualquier grupo discreto conmutativo.

De hecho, dado un tal grupo $G$, basta tomar su factorizaci\'{o}n en
producto de su parte libre y su parte de torsi\'{o}n (i.e, $\Z^n
\times \Z_{p_1^{r_1}} \times \cdots \times \Z_{p_t^{r_t}}$, con
cada $p_i$ primo y $r_i \geq 1$, para cada $1 \leq i \leq r$), y
combinar los teoremas anteriores de manera apropiada:
$$\bar{B}(\Z[G]) \Rightarrow \bar{B}(\Z[\Z] \otimes
\stackrel{n}{\cdots} \otimes \Z[\Z] \otimes \Z[\Z_{p_1^{r_1}}]
\otimes \cdots \otimes \Z[\Z_{p_t^{r_t}}]) \Rightarrow$$
$$\Rightarrow \bar{B}(\Z[\Z]) \stackrel{n}{\cdots} \bar{B}(\Z[\Z])
\otimes \bar{B}(\Z[\Z_{p_1^{r_1}}]) \otimes \cdots \otimes
\bar{B}(\Z[\Z_{p_t^{r_t}}]) \Rightarrow$$ $$\Rightarrow (E(u_1,1) \otimes
\cdots \otimes E(u_n,1) \otimes E(w_1,1) \otimes \Gamma (v_1,2)
\otimes \cdots \otimes E(w_t,1) \otimes \Gamma (v_t,2), \;d),$$
donde $d=d_{p_1^{r_1}} \otimes \cdots \otimes d_{p_t^{r_t}}).$

Originalmente, en \cite{EM53} y \cite{EM54}, este proceso se
realizaba mediante la utilizaci\'{o}n del clasificante geom\'{e}trico,
que notaremos en adelante simplemente por $\bar{W}$:
a partir del producto cartesiano $G \times G'$ de dos grupos
abelianos finitamente generados $G$ y $G'$, se establec\'{\i}a un
isomorfismo simplicial trivial entre $\bar{W}(G \times G')$ y
$\bar{W}(G) \times \bar{W}(G')$, que desembocaba en el resultado
obtenido anteriormente a partir de los modelos de los factores.

Nosotros pretendemos seguir un proceso similar en el caso de los
productos semidirectos de grupos.

Recordemos que, dados dos grupos $(A,+)$ y $(G,\cdot)$ y una acci\'{o}n
$\chi :G \times A \rightarrow A$ de $G$ en $A$ (esto es, una aplicaci\'{o}n de
modo que $\chi (e,a)=a$ y $\chi (g'g,a)= \chi (g',(\chi (g,a))$), el {\em
producto semidirecto} $A \times _\chi G$ de $A$ y $G$ consiste en el grupo
resultante de considerar en $A \times G$ la operaci\'{o}n producto
definida por la regla: $$(a,g) \cdot (a',g')=(a+ \chi
(g,a'),gg').$$ La acci\'{o}n $\chi$ se denomina {\em acci\'{o}n de grupos}
cuando verifica que $$\chi (g,a+a')= \chi (g,a) + \chi (g,a') \qquad
\forall \, g \in G, \; a,a' \in A;$$ es decir, para cada $g \in G$ fijo,
la aplicaci\'{o}n $\chi (g,-)$ resulta ser un homomorfismo de grupos.

Para simplificar la lectura
en lo que sigue, notaremos $\chi (g,a)$ simplemente por $ga$.

Un ejemplo cl\'{a}sico de producto semidirecto de grupos lo conforma
el {\em grupo dih\'{e}drico} $D_{2m}= \Z _m \times _\chi \Z _2$, para
$m \geq 2$, con acci\'{o}n determinada por $\chi (n,-1)=-n$.

La noci\'{o}n de producto semidirecto admite una generalizaci\'{o}n al
contexto de los grupos simpliciales de forma natural, considerando
acciones de grupos en el sentido simplicial.

Es m\'{a}s, teniendo en cuenta la versi\'{o}n simplicial can\'{o}nica
$^sG$ asociada a todo grupo discreto $G$ (constituida por el
propio grupo $G$ en cada grado y el homomorfismo identidad como
\'{u}nico operador de cara y de degeneraci\'{o}n), podemos aglutinar
toda la informaci\'{o}n homol\'{o}gica de un producto semidirecto de
grupos $A \times _\chi G$ en su versi\'{o}n simplicial $^s(A \times
_\chi G)=\,^sA \times _{^s \chi} \, ^sG$, donde $^s \chi$ denota la
acci\'{o}n simplicial de $^sG$ en $^sA$ inducida por $\chi$. El problema
reside, pues, en encontrar un modelo homol\'{o}gico de $\bar{W}(^s(A \times
_\chi G))$.

Nos vamos a reducir, en principio, al estudio de los productos
semidirectos de grupos $A \times _\chi G$, con $A$ y $G$ grupos
simpliciales y $\chi$ acci\'{o}n de grupos.

Consideremos el producto cartesiano torcido $\bar{W}(A) \times _
\tau \bar{W}(G)$ definido por la funci\'{o}n de torsi\'{o}n $\tau_*$,
$$\tau _n: \bar{W}_n(G) \rightarrow G_{n-1}, \qquad [g_{n-1},
\ldots ,g_0] \stackrel{\tau _n}{\mapsto} g_{n-1};$$ con grupo
estructural $G$ de acci\'{o}n $\kappa$ determinada por
$$\begin{array}{ccc} \kappa\,:\,G\times \bar{W}(A) &
\longrightarrow & \bar{W}(A) \\ (g_n,[a_{n-1},\ldots,a_0]) &
\longrightarrow & [\partial _0g_n \cdot a_{n-1},\ldots, \partial _0^ng_n
\cdot a_0].
\end{array}$$

En estas circunstancias, podemos enunciar el siguiente teorema.

\begin{thr}\label{tic}
Existe un isomorfismo simplicial $$f:\bar{W}(A\times_\chi
G)\longrightarrow \bar{W}(A)\times_\tau\bar{W}(G),$$
definido por las aplicaciones\\

%\vspace*{.5cm}
\noindent $f_n[(a_{n-1},g_{n-1}),\ldots,(a_0,g_0)]  = $\\ $\quad
([g^{-1}_{n-1} a_{n-1},\ldots,
\partial^{i-1}_0g^{-1}_{n-1}\cdots \partial_0g^{-1}_{n-i+1}
g^{-1}_{n-i} a_{n-i},\ldots,
 \partial^{n-1}_0g^{-1}_{n-1}\cdots \partial_0g^{-1}_1
g^{-1}_0a_0],$\\ \hspace*{8cm}{$\qquad \qquad \qquad
[g_{n-1},\ldots,g_0]);$}

\vspace*{.25cm}
\noindent y

\vspace*{.25cm}
\noindent
$f^{-1}_n([a_{n-1},\ldots,a_0],\,[g_{n-1},\ldots,g_0]) = $\\
$\qquad \quad [(g_{n-1}a_{n-1},g_{n-1}),\ldots,
(g_{n-i}\partial_0g_{n-i+1}\cdots\partial_0^{n-i+1}g_{n-1}a_{n-i},
g_{n-i}),$\\ \hspace*{5cm}$\qquad \qquad \qquad \ldots, (g_0\partial_0g_1\cdots
\partial_0^{n-1}g_{n-1} a_0, g_0)].$
\end{thr}

Este resultado t\'{e}cnico es la llave que permite conectar con los
estudios realizados a finales de la d\'{e}cada pasada por Lambe y
Stasheff sobre la homolog\'{\i}a de fibrados iterados en
\cite{LS87}.

Aunque el PCT $\bar{W}(A)\times_\tau\bar{W}(G)$ que define un
producto semidirecto de grupos simpliciales $A \times _\chi G$ no
es principal (la acci\'{o}n se efect\'{u}a a trav\'{e}s del grupo $G$ y
no de $\bar{W}(G)$), la demostraci\'{o}n del Teorema 3.5 de
\cite{LS87} se ajusta a nuestras circunstancias, de modo que es
posible establecer un modelo homol\'{o}gico para ciertos productos
semidirectos de grupos.

\begin{thr} Sea $A \times _\chi G$ un producto semidirecto de grupos
simpliciales de modo que la cocadena de torsi\'{o}n $t$ en
$\bar{B}(A)\otimes_t \bar{B}(G)$ inducida por $\tau$ seg\'{u}n el
teorema de Eilenberg-Zilber torcido se anule actuando sobre
elementos de grado 1. En estas circunstancias, $A \times _\chi G$
tiene un modelo homol\'{o}gico, en funci\'{o}n de sendos modelos
homol\'{o}gicos de los factores $A$ y $G$.
\end{thr}

\begin{rmk} {\em La condici\'{o}n sobre la cocadena $t$ fue
considerada ya por May en \cite{May67}, quien apunt\'{o} que grupos
simpliciales como los reducidos se encuentran entre aquellos que
verifican tal premisa.}
\end{rmk}

Esto resuelve la cuesti\'{o}n para aquellos productos semidirectos
de grupos simpliciales verificando la condici\'{o}n sobre la
cocadena $t$. El problema que se presenta ahora es que los
grupos simpliciales inducidos por los grupos discretos no
verifican tal condici\'{o}n, al estar concentrados en grado cero.

Nosotros proponemos aqu\'{\i} un camino alternativo, que facilita la
Teor\'{\i}a de Perturbaci\'{o}n Homol\'{o}gica utilizando los modelos
homol\'{o}gicos de $\Z$ y $\Z _{p^r}$ anteriormente expuestos, as\'{\i} como
el isomorfismo del Teorema \ref{tic}. El teorema principal de este
art\'{\i}culo recoge dicho resultado.

\begin{thr} \label{mhps} Sean $G$ abeliano finito, $A$ un grupo discreto con
modelo homol\'{o}gico $hA$ y $\chi :G \times A \rightarrow A$ una acci\'{o}n
de grupos. Entonces, el grupo $A \times _\chi G$ tiene un modelo
homol\'{o}gico.
\end{thr}

Aun m\'{a}s, demostraremos que la diferencial sobre el modelo anterior obtenida
por perturbaci\'{o}n se
simplifica en parte al considerar lo particular de la cocadena de
Brown generada por la funci\'{o}n de torsi\'{o}n $\tau (\chi)$,
y la propia acci\'{o}n $\chi$.

En el caso particular de que la acci\'{o}n, adem\'{a}s, act\'{u}e por
componentes (en el sentido de que $\chi (g,[1,0, \ldots , 0])=[*,0, \ldots ,0]$
y as\'{\i} sucesivamente), la diferencial admitir\'{a} una expresi\'{o}n
notablemente m\'{a}s simple. Veremos que \'{e}ste es el caso de los
grupos dih\'{e}dricos, de los que calcularemos la homolog\'{\i}a.
Esto permitir\'{a} establecer una comparaci\'{o}n de nuestro m\'{e}todo
con el ya cl\'{a}sico de la sucesi\'{o}n espectral de Lindon-Hochschild-Serre (ver
\cite{Wei}).

\section{Cuestiones relacionadas.}

En este apartado nos interesaremos por resoluciones libres
asociadas al modelo anterior y por la eventual estructura de
producto tensorial torcido \'{o} $\gamma$-producto tensorial torcido
que pueda presentar.

De hecho, los resultados anteriores se pueden prolongar hasta
conectar con $A_{\infty}$-estructuras
(\cite{Sta63a,Sta63b,Sta70,Kad80,Pro84,GS86,GL89,GLS91,JL96}) y
resoluciones (\cite{Lam91,Lam92,Lam93}), de modo que podemos
establecer peque$\n$as resoluciones libres para el caso de
productos semidirectos de grupos finitos.

La obtenci\'{o}n de resoluciones libres ``econ\'{o}micas'' de grupos abelianos
(econ\'{o}micas en el sentido de que los c\'{a}lculos homol\'{o}gicos que
requieran puedan ser efectuados ciertamente en la pr\'{a}ctica), es un
problema que ha sido tratado de forma intensa en las \'{u}ltimas d\'{e}cadas
(\cite{Eil49,EM54,Car56,CE56,Tat57,Moo76,Hue91e}. Con la aparici\'{o}n de la
Teor\'{\i}a de Perturbaci\'{o}n Homol\'{o}gica
(\cite{Shi61,Bro64,Gug72,GL89,GLS91,HK91}), otros grupos discretos
comenzaron a ser tratados. En \cite{LS87}, Lambe y Stasheff describen
un m\'{e}todo para calcular la homolog\'{\i}a de extensiones centrales de
grupos. Este m\'{e}todo ya ha sido ampliado para el caso de grupos nilpotentes
libres de torsi\'{o}n finitamente generados (\cite{LS87,Lam89,Lam91}) y
$p$-grupos finitos (\cite{Lam92,Lam,EGL97}). Utilizando un proceso de
perturbaci\'{o}n paralelo, en \cite{Hue89} y \cite{Hue91t} Huebschmann se
interesa, respectivamente, por el problema del c\'{a}lculo de la
homolog\'{\i}a de grupos nilpotentes finitamente generados de dos pasos y
grupos metac\'{\i}clicos.

M\'{a}s concretamente, en \cite{LS87} se factoriza un modelo simplicial
$\bar{W}(E)$ del espacio clasificante de una extensi\'{o}n central de grupos
$E$ en forma de {\em producto cartesiano torcido}. La aplicaci\'{o}n del
Teorema de Eilenberg-Zilber Torcido, algunas propiedades de anulaci\'{o}n de
la cocadena de torsi\'{o}n asociada y la existencia de un {\em modelo
homol\'{o}gico} del grupo de enteros $\Z$ permiten finalmente construir una
resoluci\'{o}n peque$\n$a libre para grupos nilpotentes libres de torsi\'{o}n
finitamente generados. Nuestra aproximaci\'{o}n es similar, tomando
como punto de partida el modelo homol\'{o}gico hallado en la
secci\'{o}n anterior.

\begin{thr} \cite{AAR99b}
Sea $A$ una DGA-\'{a}lgebra, $C$ una DGA-co\'{a}lgebra conexa,
$t:C \rightarrow A$ una cocadena de torsi\'{o}n y
$c:\{C,H,f,g,\phi \}$ una contracci\'{o}n que induce en $H$ una
estructura de $A_{\infty}$-co\'{a}lgebra, de modo que $t \phi =0$ y
la composici\'{o}n $(t \cap )(1 \otimes \phi)$ es puntualmente
nilpotente.

Entonces, existe una contracci\'{o}n de $A \otimes _t C$ en el
$\gamma$-producto tensorial torcido $A \otimes _{t'} H$, donde
$t'=tg:H \rightarrow A$ es la $\gamma$-cocadena inducida por $c$ y
$t$.
\end{thr}

\begin{rmk} {\em Conviene hacer las siguientes puntualizaciones:
\begin{itemize}
\item La hip\'{o}tesis de conexi\'{o}n sobre $C$ se puede reducir
a exigir simplemente que $C$ est\'{e} aumentada, de modo que tenga
sentido el contemplar el DG-m\'{o}dulo $\bar{C}=\mbox{Ker }(\chi
_{\scst C})$. En el caso particular de que $C$ sea conexo, entonces
$\bar{C}_0=0$ y $\bar{C}_i=C_i$ para $i>0$.

\item Una condici\'{o}n suficiente para que $c$ induzca
sobre $H$ una estructura de $A_{\infty}$-co\'{a}lgebra es que $C$
sea simplemente conexo (ver \cite{GS86}).

\item Por otro lado, una condici\'{o}n suficiente para que la
composici\'{o}n $t \cap (1 \otimes \phi )$ sea puntualmente
nilpotente es que $t$ se anule sobre $C_1$ (ver \cite{May67},
\cite{LS87}). Este es el caso, obviamente, cuando $C$ es
simplemente conexo.

\item La hip\'{o}tesis de que la composici\'{o}n $t \phi$ es cero
resulta una pieza clave en la demostraci\'{o}n final, no s\'{o}lo ya
por la simplificaci\'{o}n en las f\'{o}rmulas que aparecen. Estamos
convencidos de que en otro caso la tesis del resultado anterior
ser\'{\i}a falsa, aunque no disponemos de contraejemplo alguno.
\end{itemize}}
\end{rmk}

A nosotros nos interesa particularmente el estudio del caso $C=
\bar{B}(A)$ y $t= \theta$, cocadena universal, a fin de obtener
una resoluci\'{o}n del producto semidirecto de un par de grupos
discretos.

\begin{thr}
Sea $A$ una DGA-\'{a}lgebra y $c: \{ \bar{B}(A), H, f, g ,\phi\}$
una contracci\'{o}n que induce en $H$ una estructura de
$A_{\infty}$-co\'{a}lgebra, de modo que $\phi$ incrementa el grado
simplicial de $\bar{B}(A)$ en uno y la composici\'{o}n $(\theta
\cap) (1 \otimes \phi)$ es puntualmente nilpotente.

Entonces, existe una contracci\'{o}n expl\'{\i}cita de la
resoluci\'{o}n bar $B(A)$ en el $\gamma$-producto tensorial
torcido $A \otimes _{\theta '} H$, donde $\theta ':H \rightarrow
A$ es la $\gamma$-cocadena definida como $\theta '=\theta g$.
\end{thr}

En estas circunstancias, podemos deducir la existencia de
resoluciones asociadas a ciertos productos semidirectos de grupos.

\begin{thr}
Sean $A$ y $G$ dos grupos discretos abelianos finitos y $\chi :G
\times A \rightarrow A$ una acci\'{o}n distributiva. Entonces,
existe un $\gamma$-producto tensorial torcido $\Z [A \times _\chi
G] \otimes _r H$ que escinde de la contracci\'{o}n bar
est\'{a}ndar, de modo que $H$ es un modelo homol\'{o}gico de $A
\times _\chi G$ (dotado con la estructura de
$A_\infty$-co\'{a}lgebra inducida correspondiente) y $r:H
\rightarrow \Z[A \times _\chi G]$ es una $\gamma$-cocadena de
torsi\'{o}n.
\end{thr}

\begin{rmk} {\em Dos comentarios a tener en cuenta:
\begin{itemize}
\item Consideremos dos grupos discretos $A$ y $G$ con modelos
homol\'{o}gicos conocidos $hA$ y $hG$, respectivamente, y una
acci\'{o}n distributiva $\chi$ de $G$ en $A$ de modo que $A \times
_\chi G$ admita un modelo homol\'{o}gico en funci\'{o}n de los de los
factores en la forma del Teorema \ref{mhps}. La demostraci\'{o}n
realizada anteriormente indica que en estas circunstancias el
modelo homol\'{o}gico en cuesti\'{o}n tambi\'{e}n viene dotado de una
estructura de $A_\infty$-co\'{a}lgebra inducida por la contracci\'{o}n
que da dicho modelo homol\'{o}gico.
\item No ocurre otro tanto con respecto a la existencia de una
resoluci\'{o}n que escinda de la resoluci\'{o}n bar est\'{a}ndar. Por
regla general, falla que en el \'{u}ltimo paso sea puntualmente
nilpotente la composici\'{o}n $(1 \otimes 1 \otimes \phi _{\scst G}+
1 \otimes \phi _{\scst A} \otimes g_{\scst G}f_{\scst G})d_{\scst
d_{\theta \cap}}$. Este es el caso cuando $A$ \'{o} $G$ presentan
parte libre no trivial (el operador de homotop\'{\i}a $\phi
_{\scst \Z}$ llega a aumentar el grado de filtraci\'{o}n en 1).
\end{itemize}}
\end{rmk}

A partir de aqu\'{\i}, es l\'{\i}cito preguntarse acerca de la
estructura de PTT \'{o} $\gamma$-PTT que pudiera subyacer
eventualmente en el modelo homol\'{o}gico hallado previamente.

De hecho, remont\'{a}ndonos una vez m\'{a}s al art\'{\i}culo que
inspirara el que nos lleva, es posible determinar una tal
estructura en modelos
homol\'{o}gicos de extensiones centrales de grupos descritos en
\cite{LS87}.

Dada una extensi\'{o}n central de grupos $A \times _f G$, donde $f$
es el 2-cociclo de la extensi\'{o}n, se puede calcular un modelo
homol\'{o}gico suyo estableciendo, paso a paso, la siguiente cadena
de contracciones, conocidos sendos modelos homol\'{o}gicos $hA$ y
$hG$ de $A$ y $G$, respectivamente:
\begin{itemize}
\item $C_*(\bar{W}(A\times_f G))\Rightarrow C_*(\bar{W}(A)\times_\tau
\bar{W}(G))$.
\item $C_*(\bar{W}(A)\times_\tau \bar{W}(G))\Rightarrow
C_*(\bar{W}(A))\otimes_{t(\tau)}C_*(\bar{W}(G)) $.
\item $C_*(\bar{W}(A))\ot_{t(\tau)}C_*(\bar{W}(G))\Rightarrow (hA\ot hG, \;d^\ot+
d_{\scst t(\tau)\cap})$.
\end{itemize}

Demostramos que bajo ciertas
hip\'{o}tesis sobre los factores de una extensi\'{o}n central de
grupos, un modelo homol\'{o}gico de tal objeto es en realidad un
producto tensorial torcido.

\begin{thr}
Sean $A$ y $G$ dos grupos abelianos finitamente generados, $G$
libre de torsi\'{o}n, y $f$ un 2-cociclo que define una extensi\'{o}n
de $G$ por $A$. Entonces, el complejo $(hA\ot hG, \;d^\ot+
d_{\scst t(\tau)\cap})$ tiene estructura de producto tensorial
torcido.
\end{thr}

\'{E}ste no es el caso, en general, del modelo homol\'{o}gico de productos
semidirectos de grupos discretos finitos que describimos en el ep\'{\i}grafe
anterior.

Sin embargo, s\'{\i} es posible determinar ciertas estructuras de
$\gamma$-producto tensorial torcido asociadas, en un sentido que
precisaremos en la versi\'{o}n final del art\'{\i}culo.

\begin{thebibliography}{99}

\bibitem{AAR98} \'{A}lvarez V., Armario J.A. y Real P.: {\em On the
homology of semidirect product of groups}, Colloquium on Topology,
Gyula (Hungr\'{\i}a) (1998).

\bibitem{AAR99} \'{A}lvarez V., Armario J.A. y Real P.: {\em
Estructuras algebraicas en Topolog\'{\i}a Algebraica y \'{A}lgebra
Homol\'{o}gica}. En preparaci\'{o}n.

\bibitem{AAR99b} \'{A}lvarez V., Armario J.A. y Real P.: {\em On
resolutions which split off of the Bar construction}. En
preparaci\'{o}n.

\bibitem{Bro82} Brown K.S.: {\em Cohomology of Groups}, Springer,
Berlin (1982).

\bibitem{Bro59} Brown E.H.: {\em Twisted tensor products I},
Annals of Math., {\bf 69}, pp. 223-246 (1959).

\bibitem{Bro64} Brown R.: {\em The twisted Eilenberg-Zilber theorem},
Celebrazioni Achimedee del secolo XX, Simposio di topologia
(Messina, 1964). Ed. Oderisi, Gubbio, pp. 33-37 (1965).

\bibitem{Car56} Cartan H.: {\em Alg\`ebres d'Eilenberg-Mac Lane},
S\'{e}minaire H. Cartan 1954/55, (expos\'{e} 2 \`a 11), Ecole Normale
Superi\`ere, Paris (1956).

\bibitem{CE56} Cartan H. y Eilenberg S.: {\em Homological Algebra},
Princeton University Press, Princeton (1956).

\bibitem{Eil49} Eilenberg S.: {\em Topological methods in abstract
algebra}, Bull. Amer. Math. Soc., {\bf 55}. pp. 3-27 (1949).

\bibitem{EM53} Eilenberg S. y Mac Lane S.: {\em On the groups
$H(\pi, n)$, I}, Annals of Math., {\bf 58}, pp. 55-106 (1953).

\bibitem{EM54} Eilenberg S. y Mac Lane S.: {\em On the groups
$H(\pi, n)$, II}, Annals of Math., {\bf 66}, pp. 49-139 (1954).

\bibitem{EZ59} Eilenberg S. y Zilber J.A.: {\em On the products
of complexes}, Am. J. Math., {\bf 75}, pp. 200-204 (1959).

\bibitem{EGL97} Ekedahl T., Grabmeier J. y Lambe L.: {\em Algorithms
for algebraic computations with applications to the cohomology of
finite $p$-groups}, Preprint of Department of Math. and Center
for Innovative Computation, University of Wales (1997).

\bibitem{Gug60} Gugenheim V.K.A.M.: {\em On a theorem of E.H.
Brown}, IL. J. Math., {\bf 4}, pp. 292-311 (1960).

\bibitem{Gug72} Gugenheim V.K.A.M.: {\em On a the chain complex of a
fibration}, Illinois J. Math., {\bf 3}, pp.398-414 (1972).

\bibitem{GL89} Gugenheim V.K.A.M. y Lambe L.: {\em Perturbation theory
in Di\-ffe\-ren\-tial Homological Algebra, I}, Illinois J. Math., {\bf
33}, pp. 556-582 (1989).

\bibitem{GLS90} Gugenheim V.K.A.M., Lambe L. y Stasheff J.: {\em
Algebraic aspects of Chen's twisting cochain}, Illinois J. Math.,
{\bf 34}, pp. 485-502 (1990).

\bibitem{GLS91} V.K.A.M. Gugenheim, L. Lambe and
J. Stasheff, Perturbation theory in Di\-ffe\-rential Homological
Algebra  II, {\em Illinois J. Math. \/} {\bf 35}
n. 3 (1991) 357--373.

\bibitem{GM74MAMS} Gugenheim V.K.A.M. y May J.P.: {\em On the theory
and application of differential torsion products}, Memo. Amer. Math. Soc.,
{\bf 142} (1974).

\bibitem{GS86} Gugenheim V.K.A.M. y Stasheff J.: {\em On
perturbations and $A_{\infty}$ structures}, Bull. Soc. Math.
de Belg., {\bf 38}, pp. 237-246 (1986).

\bibitem{Hue89} Huebschmann J.: {\em Cohomology of nilpotent groups of
class 2}, J. Alg., {\bf 126}, pp. 400-450 (1989).

\bibitem{Hue91e} Huebschmann J.: {\em Cohomology of finitely
generated abelian groups}, L'Enseignement Math\'{e}matique {\bf 37},
pp. 61-71 (1991).

\bibitem{Hue91t} Huebschmann J.: {\em Cohomology of metacyclic groups},
Transactions of the American Mathematical Society, {\bf  328}, {\bf n. 1},
pp. 1-72 (1991).

\bibitem{HK91} Huesbschmann J. y Kadeishvili T.K.: {\em Small models
for chain algebras}, Math. Z., {\bf 207}, pp. 245-280 (1991).

\bibitem{JL96} Johansson L. y Lambe L.: {\em Transferring algebra
structures up to homology equivalence}, Math. Scand (1996).

\bibitem{Kad80} Kadeishvili T.K.: {\em On the homology theory of
fibre spaces}, Russian Math. Surveys, {\bf 35}, pp. 231-238
(1980).

\bibitem{Lam89} Lambe L.: {\em Algorithms for the homology of
nilpotent groups}, Conf. on applications of computers to Geom.
and Top., Lecture Notes in Pure and Applied Math.,
Marcel Dekker Inc., N.Y., {\bf 114} (1989).

\bibitem{Lam91} Lambe L.: {\em Resolutions via homological
perturbation}, J. Symbolic Comp., {\bf 12}, pp. 71-87 (1991).

\bibitem{Lam92} Lambe L.: {\em Homological perturbation theory,
Hochschild homology and formal groups}, Proc. Conference on Deformation
Theory and Quantization with Applications to Physics, Amherst, MA, June
1990, Cont. Math., {\bf 134} A.M.S, pp. 183-218 (1992).

\bibitem{Lam93} Lambe L.: {\em
Resolutions which split off of the bar
construction}, Journal of Pure and Applied Algebra, {\bf 84}, pp.
311-329 (1993).

\bibitem{Lam} Lambe L.: {\em An algorithm for calculating cocycles},
Preprint of Department of Mathematics and Center for
Innovative Computation, University of Wales (1997).

\bibitem{LS87} Lambe  L. y Stasheff J.: {\em Applications of
perturbation theory to iterated fibrations}, Manuscripta Math.,
{\bf 58}, pp. 367--376 (1987).

\bibitem{Mac95} Mac Lane S.: {\em Homology}, Classics in
Mathematics, Springer-Verlag, Berlin (1995). Reedici\'{o}n del original
de 1975.

\bibitem{May67} May J.P.: {\em Simplicial objects in Algebraic Topology},
Chicago Lectures in Mathematics, The University of Chicago
Press, Chicago and London (1992).

\bibitem{Moo76} Moore J.C.: {\em Cartan's constructions}, Colloque
analyse et topologie, en l'honneur de H. Cartan, Ast\'{e}risque,
{\bf 32-33}, pp. 173--221 (1976).

\bibitem{Pro84} Proute A.: {\em
Alg\'{e}bres defferentielles fortement homotopiquement associatives},
 Th\`{e}se de Mathematiques, Universit\'{e} Paris VII (1984).

\bibitem{Rea95} Real P.: {\em Algoritmos de c\'{a}lculo de
homolog\'{\i}a efectiva de los espacios clasificantes}, Tesis,
Universidad de Sevilla (1993).

\bibitem{Rea96} Real P.: {\em Homological Perturbation
Theory and Associativity}, Preprint del Dpto. de Matem\'{a}tica
Aplicada I de la Universidad de Sevilla (1996). Enviado a Journal
of Pure and Applied Algebra.

\bibitem{RS94} Real P. y Sergeraert F.: {\em A conjecture of
Eilenberg-Mac Lane}, Preprint de l'Institute Fourier, Grenoble
(1995).

\bibitem{Shi61} Shih W.: {\em Homologie des espaces fibr\'{e}s}, Inst.
Hautes Etudes Sci., {\bf 13}, pp. 293-312 (1962).

\bibitem{Sta63a} Stasheff J.: {\em Homotopy associativity of a $H$-space
I}, Trans. Amer. Math. Soc., {\bf 108}, pp. 275-292 (1963).

\bibitem{Sta63b} Stasheff J.: {\em Homotopy associativity of a
$H$-space II}, Trans. Amer. Math. Soc., {\bf 108}, pp. 293-312
(1963).

\bibitem{Sta70} Stasheff J.: {\em $H$-spaces from a homotopy point of
view}, Lecture Notes in Math., {\bf 161}, Springer, N.Y., (1970).

\bibitem{Tat57} Tate J.: {\em Homology of Noetherian and local rings},
Illinois J. of Math., {\bf 1}, pp. 14-27 (1957).

\bibitem{Wei} Weibel C.A.: {\em An introduction to homological
algebra}, Cambridge studies in advanced mathematics {\bf 38}, Cambridge
University Press (1994).

\end{thebibliography}


\end{document}

